\section{B570 Repair, Missingness, Suffixes, and Paid Costs}
\label{supp:b570}
\subsection{Earlier finite closeout: a positive local case and failures}
The earlier B570 closeout's normative audit establishes completion of 1,456 core records (\texttt{COMPLETE\_FINITE\_SCOPE\_AUDITED}). Historical report wording about incomplete normative closure is superseded by that completed audit, not treated as a still-open experiment. The original findings remain finite: on the shared 24-task MBPP panel, the repaired old N1 qualifies but gives only $0.188525\times/0.179055\times$ generation/request speed against ordinary $R$, and $0.804023\times/0.693541\times$ against matched target-only. The old N0/N1 comparison is $0.999522\times/1.247059\times$: request progress against N0 does not imply ordinary-reference benefit.

On eight project-unused MBPP inputs, N1 and its matched path fail MBPP/72 in all four repeats. Five UB2 alternatives retain failures involving HumanEval/42, /145 or /138 and receive a NO-GO decision. The local $k=3$/UB4/cap128 candidate on seven exposed longer probes is a positive finite case, with $1.345798\times$ generation and $1.048339\times$ request speed. It does not replace the failed expanded qualification. Long-cap anchor qualification passes 84/84 at caps256/512, with longest actual response 261 IDs. No natural state cross occurs in this earlier study; the crosses in Section~\ref{supp:diagnosis} are later experiments.

A wrongly placed UHD730 campaign of 190 requests is rejected as infrastructure evidence rather than scientific B570 data. Its records and recovery remain preserved. Recovery does not authorize outcome-dependent retuning or replacement of failed tasks. The later phase-binding candidate is a distinct implementation.

\subsection{Same-context phase binding}
The B570 candidate preserves one live target context and sampler. Ordinary-reference prefill uses UB128 with caller output reservation 512/per-sequence reservation 1; the frontend samples and accepts the pending first token once. It then binds continuation UB1 and synchronizes/reserves the scheduler before native $k=2$, draft-UB1 proposal, acceptance, bonus and rollback. The private API patch changes phase parameters and reservation within the context; it adds no model arithmetic, Vulkan shader or precision change. It has no task-ID routing, known reference token, margin-based correction or extra reference forward.

Unlike old N1, it removes a second target context, state handoff and three serialized witness readbacks. Capture-on checks compare KV before and after binding; capture-off avoids witness readback. Those checks do not supply a separate profile of the saved cost. The frontend still pays for synchronization, phase reservation, draft work, target work and reconciliation. S2 includes \texttt{phase.h}, \texttt{target.cpp}, \texttt{native.cpp} and the private-library delta with source bindings.

\subsection{Qualification counts and unresolved raw comparisons}
Table~\ref{tab:b570_qualification} retains endpoint and raw-array distinctions. The 48 baseline-development rows have only 32 ordinary-reference-equal endpoints; they are not all qualified candidate evidence. The 12 shams are separate from the 299 rows in the qualification export. Stage D contains 56 exposed tasks evaluated twice. Stage E combines eight bridge records and 119 formal records; their roles are not interchangeable.

\begin{table}[ht]\centering\small
\caption{B570 candidate qualification stages and raw comparison limits. Equality is to the registered endpoint reference unless stated otherwise.}
\label{tab:b570_qualification}
\begin{tabularx}{\linewidth}{@{}lrrX@{}}\toprule
Stage & Candidate checks & Endpoint equal & Raw reference comparison\\\midrule
C development & 12 & 12 & Four equal, eight different\\
D exposed bridge & 112 & 112 & MBPP: 16 equal/48 different; 48 HumanEval raw comparisons unknown\\
E candidate formal & 30 & 30 & Included in the E total below; suffix cost retained\\
E all available plus bridge & 127 & 127 & 80 equal, 47 different\\\bottomrule
\end{tabularx}
\end{table}

The four-condition formal matrix is 119/120. The missing unit is \nolinkurl{E/FORMAL/1/old_N1/MBPP/400}, with \nolinkurl{INFRASTRUCTURE_INTERRUPTION_NO_PAYLOAD}. Its runtime cost and interrupted teardown time remain unknown. Each of ordinary $R$, matched $M$ and candidate $C$ has 30 complete formal observations and 15 complete two-round task pairs. Old N1 has 29 observations and only 14 complete task pairs. The available MBPP/400 old-N1 round is retained but cannot supply its two-round median. There is no interpolation, zero-cost assumption, silent task removal from other comparisons or scientific rerun. The study's administrative closeout remains \nolinkurl{PARTIAL_INFRASTRUCTURE_BUDGET_CLOSED}.

\subsection{All measured timers and work}
Generation time is the frontend's registered C++ interval, including pending-state initialization, sampler/batch operations, phase binding within that interval, draft/verification, rollback and final synchronization. Whole-request time additionally includes prefill/reset/cleanup, excluding subsequent disk file writes. Host request time is separately measured by Windows QPC. Earlier A/B stage GetTickCount's 15.625-ms granularity is not substituted for formal C/D/E QPC measurements. Model-load cost is outside these comparisons.

Table 4 of the main manuscript reports all three formal timer aggregates. Table~\ref{tab:b570_task_cost} provides all paired task generation and whole-request ratios. R/C and M/C are below one on every complete task and round. Old-N1/C whole and host request ratios reverse on MBPP/245 in round 2, approximately 0.873582 and 0.873603; aggregate old-path improvement therefore does not imply universal improvement.

% Generated by scripts/reconstruct.py from frozen exports.
\begingroup\small
\begin{longtable}{@{}rrrrrrr@{}}
\caption{All B570 two-round task-median speed ratios. G/W denote generation/whole request; old N1 has no complete pair for MBPP/400.}\label{tab:b570_task_cost}\\
\toprule
MBPP & $R/C$ G & $R/C$ W & $M/C$ G & $M/C$ W & N1/$C$ G & N1/$C$ W\\\midrule
\endfirsthead
\toprule
MBPP & $R/C$ G & $R/C$ W & $M/C$ G & $M/C$ W & N1/$C$ G & N1/$C$ W\\\midrule
\endhead
\midrule\multicolumn{7}{r}{Continued on next page}\\
\endfoot
\bottomrule
\endlastfoot
125 & 0.320296 & 0.270051 & 0.319891 & 0.243047 & 2.437945 & 1.908679\\
130 & 0.299904 & 0.252594 & 0.304551 & 0.257891 & 2.269168 & 1.978931\\
132 & 0.264679 & 0.162455 & 0.271252 & 0.169635 & 8.515349 & 4.089033\\
245 & 0.298694 & 0.254655 & 0.298448 & 0.255268 & 1.541642 & 1.376354\\
278 & 0.270026 & 0.205714 & 0.271973 & 0.245382 & 3.078819 & 2.321202\\
312 & 0.257757 & 0.195314 & 0.259093 & 0.198785 & 4.282380 & 2.984986\\
400 & 0.265906 & 0.153521 & 0.269833 & 0.158972 & -- & --\\
429 & 0.276956 & 0.197973 & 0.281509 & 0.203527 & 3.441493 & 2.411608\\
435 & 0.260099 & 0.151096 & 0.266645 & 0.158111 & 6.873734 & 3.323427\\
446 & 0.305206 & 0.233673 & 0.306943 & 0.272100 & 2.479117 & 1.929448\\
451 & 0.273297 & 0.193117 & 0.276409 & 0.198226 & 4.774253 & 3.079460\\
470 & 0.281105 & 0.243139 & 0.284055 & 0.199636 & 4.389394 & 2.870175\\
475 & 0.276814 & 0.196772 & 0.278850 & 0.201003 & 5.011890 & 3.193580\\
479 & 0.310439 & 0.196746 & 0.312663 & 0.200719 & 5.061865 & 2.965935\\
69 & 0.286855 & 0.163151 & 0.292541 & 0.169808 & 8.776569 & 4.056630\\
\end{longtable}

\endgroup


Candidate formal endpoint IDs total 1,150 and raw IDs 1,176. All 26 suffix IDs are charged. Drafted/accepted IDs total 888/702, and target/draft continuation caller tokens 1,332/1,362. These denominators describe different runtime quantities and cannot be collapsed into one acceptance-adjusted cost. The narrowed target common-sampling measurement of 851.52 ms omits other work; it is not a full component profile. The known complete-campaign host subtotal is approximately 6,541.86 s, excluding the unknown interruption/teardown. It is not the cost of every planned request or a normalized per-task speed estimator.
